\documentclass[10pt,a4paper]{amsart}
\usepackage[margin=19mm]{geometry}
\usepackage{amsmath,amssymb,mathtools}
\usepackage[scr=rsfs,cal=euler]{mathalfa}
\usepackage{xfrac}
\usepackage[unicode,hidelinks,bookmarks=false]{hyperref}
\newcommand{\ie}{i.e.\ }
\newcommand{\cf}{cf.\ }
\newcommand{\resp}{respectively\ }
\newcommand{\Subsection}{\subsection}
\providecommand{\nobreakdash}{\nobreak\hbox{-}}
\setlength{\emergencystretch}{2em}
\multlinegap0pt
\usepackage{bm}
\usepackage{bbm}
\usepackage{dsfont}
\usepackage{xspace}
\usepackage{cancel}
\usepackage{mathtools}
\usepackage{cleveref}
\usepackage{amsthm}
\makeatletter
\@ifundefined{theorem}{\newtheorem{theorem}{Theorem}}{}
\@ifundefined{lemma}{\newtheorem{lemma}[theorem]{Lemma}}{}
\@ifundefined{remark}{\newtheorem{remark}[theorem]{Remark}}{}
\makeatother
\DeclareMathOperator{\Dom}{Dom} % definition domain
\title[An equivariant compactification for adjoint reductive group ]{An equivariant compactification for adjoint reductive group schemes}
\author{Shang Li}
\date{Source: arXiv:2308.01715v5 (CC BY 4.0)}
\begin{document}
\maketitle

\noindent\emph{Source and licence.} An equivariant compactification for adjoint reductive group schemes. Version arXiv:2308.01715v5 (CC BY 4.0), distributed under the Creative Commons Attribution 4.0 International licence, as recorded in the arXiv licence metadata for this exact version and shown on the abstract page https://arxiv.org/abs/2308.01715v5. Full rights notice: licenses/arxiv-2308.01715-CC-BY-4.0.txt.\par\smallskip

\noindent\textit{Editorial scope.} Counted unit: the Remark whose source label is \texttt{None} in the article (source0.tex lines 463--465, source label \texttt{None}), delivered together with the article's own complete original proof of it (source0.tex lines 467--532, one \texttt{proof} environment of 66 lines). No mathematics is written, repaired or re-derived: statement and proof are the article's own text line for line. Required context retained uncounted: Omegaembedformula (lines 301--303); extensionofmultiplication (lines 386--390); theoremrationalaction (lines 451--461). Every definition, display and label of those lines is retained unchanged. Omitted from this package and not redistributed: 1--300, 304--385, 391--450, 462--462, 466--466, 533--1379. Nothing inside a retained excerpt is omitted or altered: every retained body is delivered exactly as the article's lines are written, so no omission receipt of any kind accompanies this package. Citations: the retained excerpts carry no citation command at all, so no bibliography entry of the article is transcribed and no reference is redistributed. Reference adaptation: none was needed and none was made; every reference inside the delivered unit resolves inside it, so no unresolved reference or citation is produced. Printed slips kept literal: no printed word, symbol, hypothesis or number of the counted statement or of its proof was altered. Layout and class: the article's own class is \texttt{amsart} with packages including geometry, fontenc, inputenc, color, babel, biblatex, amsmath, amsthm, mathabx, amscd, xy, hyperref, tikz-cd, enumitem; the wrapper is the lane's pinned \texttt{amsart} editorial wrapper at 10pt on A4, which restates the theorem-like environments the retained text needs and adds the article's own macro definitions that the retained text needs (\texttt{\textbackslash Dom}), with extra wrapper packages (bm, bbm, dsfont, xspace, cancel, mathtools, cleveref). The delivered numbering is the unit's own. Assets: no retained span contains a figure environment, a graphics-inclusion command, a PDF-page inclusion, a table or a TikZ picture; no image file is copied into this package, the package declares no asset dependency and no image file is redistributed with it. Licence: version arXiv:2308.01715v5 of this article is distributed under the Creative Commons Attribution 4.0 International licence, as recorded in the arXiv licence metadata for this exact version and shown on the abstract page https://arxiv.org/abs/2308.01715v5. A full rights notice is delivered at licenses/arxiv-2308.01715-CC-BY-4.0.txt. Characters: every retained excerpt is pure ASCII, so the delivered engine, XeLaTeX with its default Latin Modern text font, reports no missing character.
\begin{equation}\label{Omegaembedformula}
    \nu\colon\Omega_G:=U^-\times T\times U^+\longrightarrow \overline{\Omega}:=U^-\times \overline{T}\times U^+
\end{equation}

\begin{lemma}\label{extensionofmultiplication}
    Consider the $S$-rational morphism 
    $$\Theta: U^+\times \overline{T}\times U^-\dashrightarrow \overline{\Omega}$$
    given by the group multiplication $U^+\times T\times U^-\longrightarrow G$ and the open immersion $\nu$ \Cref{Omegaembedformula} where $\Omega_G$ is identified with the big cell of $G$ via the group multiplication. Then $\Dom(\Theta)$ contains $e\times \overline{T}\times e$.
\end{lemma}

\begin{theorem}\label{theoremrationalaction}
    There exists a unique $S$-rational morphism $\pi\colon G\times\overline{\Omega}\times G\dashrightarrow \overline{\Omega}$ whose definition domain $\Dom(\pi)$ contains $e\times\overline{\Omega}\times e$ 
    such that 

    (1) The restriction of $\pi$ to the open subscheme $(G\times\Omega_G\times G)\bigcap\Dom(\pi)$ is given by 
$$(g_1,g,g_2)\longmapsto g_1gg_2^{-1}.$$
        
    (2) For every $x\in S$, the geometric fiber $\pi_{\overline{k(x)}}$ agrees with the restriction of the rational action of $G_{\overline{k(x)}}\times G_{\overline{k(x)}}$ on the big cell $\overline{\Omega}_{\overline{k(x)}}$ of the wonderful compactification of $G_{\overline{k(x)}}$ over the algebraically closed field $\overline{k(x)}$.

    (3) The restriction $\pi\vert_{e\times\overline{\Omega}\times e}$ is the projection onto the middle factor. 
\end{theorem}

\begin{remark}
    In the case when $S$ is the spectrum of an algebraically closed field, if we assume the existence of the classical wonderful compactification, the desired $S$-rational morphism $\pi$ in \Cref{theoremrationalaction} exists simply as a restriction of the group action morphism of the wonderful compactification of $G$.  
\end{remark}

\begin{proof}
    Note that (2) follows from (1) by the $S$-density of $\Omega_G$ in $\overline{\Omega}$. The \emph{uniqueness} of $\pi$ also follows from this reason and (1).

    
    For the \emph{existence} of $\pi$, since $\Omega_G$ is $S$-dense in $G$, it suffices to define an $S$-rational morphism $$\pi\colon\Omega_G\times\overline{\Omega}\times\Omega_G\dashrightarrow \overline{\Omega}$$
    which is defined over an open subscheme $\mathcal{R}$ containing $e\times\overline{\Omega}\times e$.
    We will define our $\pi$ step by step and shrink $\Omega_G\times\overline{\Omega}\times\Omega_G$ along the way to obtain the sought $\mathcal{R}$. 
    Let $S'$ be a test $S$-scheme.
    Consider an element
    $$A:=((u_1^-,t_1, u_1^+),(u^-,t, u^+),(u_2^-,t_2, u_2^+))\in (\Omega_G\times\overline{\Omega}\times\Omega_G)(S').$$
    
    %The principle behind our construction is that we need to ``move'' all components of $A$ on $U^-$ (resp., $U^+$, $\overline{T}$) together to form those of $\pi(A)$.
    
    As the \emph{first step}, we let 
    $$V\coloneq\sigma^{-1}(\Omega_G)\bigcap \Omega_G$$
    where $\sigma$ is the inverse morphism of the group scheme $G$. Then we define an $S$-dense open subscheme 
    $$\mathcal{R}_1\coloneq\Omega_G\times\overline{\Omega}\times  V\subset\Omega_G\times\overline{\Omega}\times\Omega_G.$$
    If $A\in \mathcal{R}_1(S')$, we write 
    $$\sigma(u_2^-,t_2, u_2^+)=(\hat{u}_2^-,\hat{t}_2, \hat{u}_2^+)\in \Omega_G(S').$$
    
    As the \emph{second step}, we define 
    the $S$-dense open subscheme $\mathcal{R}_2\subset\mathcal{R}_1$ given by the following condition:
    $$A\in \mathcal{R}_2\iff
    u_1^+u^-\, \text{and}\,u^+\hat{u}_2^-\, \text{lie in}\; \Omega_G(S').$$
    If so, we write 
    \begin{equation}\label{comutequation1}
        u_1^+u^-=(\dot{u}_1^-, \dot{t}_1, \dot{u}^+)\in\Omega_G(S')
    \end{equation} 
    and 
    \begin{equation}\label{comutequation2}
        u^+\hat{u}_2^-=(\dot{u}^- ,\dot{t}_2 ,\dot{u}_2^+)\in\Omega_G(S').
    \end{equation}
    
    
    
    %we seek to ``exchange'' the positions of the positive part and the negative part of the triple $(\dot{u}^+, t_0,\dot{u}^-)$. To simplify the notations, we do this for the triple $(\dot{t}_1\dot{u}^+\dot{t}_1^{-1},\dot{t}_1 t\dot{t}_2,\dot{t}_2^{-1}\dot{u}^-\dot{t}_2)$. To do this,  

   As the \emph{third step}, we further restrict to the open subscheme $\mathcal{R}\subset\mathcal{R}_2$ defined by the following:
    $$A\in \mathcal{R}\iff(\dot{t}_1\dot{u}^+\dot{t}_1^{-1},\dot{t}_1 t_0\dot{t}_2, \dot{t}_2^{-1}\dot{u}^-\dot{t}_2)\in \Dom(\Theta)(S')$$
    where $\Theta$ is defined in \Cref{extensionofmultiplication}. The open subscheme $\mathcal{R}$ is $S$-dense because, by \Cref{extensionofmultiplication}, $e\times \overline{T}\times e\subset \Dom(\Theta)$.
    If $A\in\mathcal{R}(S')$, we write 
    \begin{equation}\label{comutequation3}    \Theta(\dot{t}_1\dot{u}^+\dot{t}_1^{-1},\dot{t}_1 t\dot{t}_2,\dot{t}_2^{-1}\dot{u}^-\dot{t}_2)=(\Ddot{u}^-, \Ddot{t}, \Ddot{u}^+)\in\overline{\Omega}(S').
    \end{equation}
    Now we have all components of $\pi(A)$ on $U^-$, $U^+$ and $\overline{T}$ in the ``right'' places. Hence 
    we define the image of $A$ under $\pi$ to be 
    \begin{equation}\label{definitionofpi}
        (u_1^- (t_1\Dot{u}_1^- t_1^{-1})(t_1\Ddot{u}^-t_1^{-1}),t_1\Ddot{t} \hat{t}_2,(\hat{t}_2^{-1}\Ddot{u}^+\hat{t}_2)(\hat{t}_2^{-1}\dot{u}_2^+\hat{t}_2)\hat{u}_2^+)\in \overline{\Omega}(S').
    \end{equation}
     Note that, when $(u^-_1, t_1,u^+_1)=e$ and $(u^-_2, t_2,u^+_2)=e$, we have that $\dot{u}^+=e$ and $\dot{u}^-=e$. Then, by \Cref{extensionofmultiplication} (1), we conclude that $\mathcal{R}$ contains $e\times \overline{\Omega}\times e$. The claim (4) follows from the definition of $\pi$ and $\mathcal{R}$ above and \Cref{extensionofmultiplication} (2).

    To show the claim (1), according to the definition of $\pi$, if $A\in(\Omega_G\times\Omega_G\times\Omega_G)(S')\bigcap\mathcal{R}(S')$, we have that
    \begin{equation}\label{comutequation4}
        \pi(A)=u_1^-t_1\dot{u}_1^-\Ddot{u}\Ddot{t}\Ddot{u}^+\dot{u}_2^+\hat{t}_2\hat{u}_2^+
    \end{equation}
    holds in $G(S')$.
    Notice that \Cref{comutequation1} -- \eqref{comutequation3} give rise to 
    $$u_1^+u^-=\dot{u}_1^-\dot{t}_1\dot{u}^+,\;\;u^+\hat{u}_2^-=\dot{u}^-\dot{t}_2 \dot{u}_2^+\;\;$$
    $$\;\;\dot{t}_1\dot{u}^+\dot{t}_1^{-1}\dot{t}_1t\dot{t}_2\dot{t}_2^{-1}\dot{u}^-\dot{t}_2=\Ddot{u}^- \Ddot{t}\Ddot{u}^+$$
    in $G(S')$,
    where \Cref{extensionofmultiplication} is used to deduce the last equation. 
    Also since $\sigma$ is the inverse operation, combining with the first step, we have 
    $$\hat{u}_2^-\hat{t}_2\hat{u}_2^+=(u_2^-t_2u_2^+)^{-1}.$$
    Substituting the above four formulas into  \Cref{comutequation4} in turn, after a computation, we have 
    $$\pi(A)=(u_1^-t_1 u_1^+)(u^-tu^+)(u_2^-t_2u_2^+)^{-1},$$
    as desired. The claim (3) follows from the definition of $\pi$.
\end{proof}

\end{document}
